% December 236, 2002 (evening edition)
%\documentclass[12pt,reqno]{amsart}
\documentclass[12pt,reqno,a4]{article}

\usepackage{amsmath}
\usepackage{amssymb}
\usepackage{amscd}


%%%% DELETE THESE LINES FOR FINAL VERSION
\newcommand{\ver}{{\it 1}}%<---PUT VERSION NUMBER HERE
\newcommand{\club}{\clubsuit} % THIS IS TO MARK COMMENTS, CHANGES
\newcommand{\cs}{$\club$}
%\usepackage{showkeys}  % THIS SHOWS LABELS
%%%%

%\topmargin -10mm
%\textheight 230mm
%\oddsidemargin -7mm
%\evensidemargin -7mm
%\textwidth 170mm


%\topmargin -300mm  %%
%%\topmargin 10mm    %% FOR FORTH FLOOR
%\textheight 230mm
%\oddsidemargin -2mm
%\evensidemargin -2mm
%\textwidth 165mm

%\topmargin=0.6in   %%%%%%%%%%%%%%(Tatsuya: this is for my machine)

\usepackage{latexsym}
\usepackage{graphicx}
%\usepackage{eclepsf}   %%%%%%%%%% This is for figures (Tatsuya)


\newcommand{\T}{{\mathbf T}^m}
\newcommand{\boxb}{\raisebox{-.8pt}{$\Box$}_b}
\newcommand{\sm}{\setminus}
\newcommand{\szego}{Szeg\"o }
\newcommand{\nhat}{\raisebox{2pt}{$\wh{\ }$}}
\newcommand{\Si}{\Sigma}
\newcommand{\inv}{^{-1}}
\newcommand{\kahler}{K\"ahler }
\newcommand{\sqrtn}{\sqrt{N}}
\newcommand{\wt}{\widetilde}
\newcommand{\wh}{\widehat}
\newcommand{\PP}{{\mathbb P}}
\newcommand{\N}{{\mathbb N}}
\newcommand{\R}{{\mathbb R}}
\newcommand{\C}{{\mathbb C}}
\newcommand{\Q}{{\mathbb Q}}
\newcommand{\Z}{{\mathbb Z}}
\newcommand{\B}{{\mathbb B}}
\newcommand{\CP}{\C\PP}
\renewcommand{\d}{\partial}
\newcommand{\dbar}{\bar\partial}
\newcommand{\ddbar}{\partial\dbar}
\newcommand{\U}{{\rm U}}
\newcommand{\D}{{\mathbf D}}
\newcommand{\E}{{\mathbf E}}
\renewcommand{\H}{{\mathbf H}}
\newcommand{\zb}{{\mathbf Z}}
\newcommand{\half}{{\frac{1}{2}}}
\newcommand{\vol}{{\operatorname{Vol}}}
\newcommand{\diag}{{\operatorname{diag}}}
\newcommand{\codim}{{\operatorname{codim\,}}}
\newcommand{\SU}{{\operatorname{SU}}}
\newcommand{\FS}{{{\operatorname{FS}}}}
\newcommand{\supp}{{\operatorname{Supp\,}}}
\renewcommand{\phi}{\varphi}
\newcommand{\eqd}{\buildrel {\operatorname{def}}\over =}
\newcommand{\ispa}[1]{\langle \,#1 \,\rangle } %%%%%%%%%%%%% Tatsuya added
\newcommand{\tf}[2]{{\textstyle \frac{#1}{#2}}} %%%%%%%%%%%%% Tatsuya added
\newcommand{\comp}{{\scriptstyle \circ}} %%%%%%%%%%%%%%%%%%% Tatsuya added
\newcommand{\op}{{\rm Op}}
\newcommand{\mb}[1]{\mathbb{#1}}
\newcommand{\mf}[1]{\mathfrak{#1}}
\newcommand{\go}{\mathfrak}
\newcommand{\bcal}{\mathcal{B}}
\newcommand{\ccal}{\mathcal{C}}
\newcommand{\dcal}{\mathcal{D}}
\newcommand{\ecal}{\mathcal{E}}
\newcommand{\fcal}{\mathcal{F}}
\newcommand{\gcal}{\mathcal{G}}
\newcommand{\hcal}{\mathcal{H}}
\newcommand{\jcal}{\mathcal{J}}
\newcommand{\lcal}{\mathcal{L}}
\newcommand{\mcal}{\mathcal{M}}
\newcommand{\ncal}{\mathcal{N}}
\newcommand{\ocal}{\mathcal{O}}
\newcommand{\pcal}{\mathcal{P}}
\newcommand{\qcal}{\mathcal{Q}}
\newcommand{\rcal}{\mathcal{R}}
\newcommand{\scal}{\mathcal{S}}
\newcommand{\tcal}{\mathcal{T}}
\newcommand{\ucal}{\mathcal{U}}
\newcommand{\vcal}{\mathcal{V}}
\newcommand{\zcal}{\mathcal{Z}}
\newcommand{\al}{\alpha}
\newcommand{\be}{\beta}
\newcommand{\ga}{\gamma}
\newcommand{\Ga}{\Gamma}
\newcommand{\La}{\Lambda}
\newcommand{\la}{\lambda}
\newcommand{\ep}{\varepsilon}
\newcommand{\de}{\delta}
\newcommand{\De}{\Delta}
\newcommand{\om}{\omega}
\newcommand{\Om}{\Omega}
\newcommand{\di}{\displaystyle}
\newtheorem{theo}{{\sc Theorem}}[section]
\newtheorem{maintheo}{{\sc {\bf Theorem}}}
\newtheorem{cor}[maintheo]{{\sc Corollary}}
\newtheorem{maincor}[maintheo]{{\sc Corollary}}
\newtheorem{lem}{{\sc Lemma}}
\newtheorem{prop}{{\sc Proposition}}
\newenvironment{example}{\medskip\noindent{\it Example:\/} }{\medskip}
\newenvironment{rem}{\medskip\noindent{\it Remark:\/} }{\medskip}
\newenvironment{defin}{\medskip\noindent{\it Definition:\/} }{\medskip}
\newenvironment{claim}{\medskip\noindent{\it Claim:\/} }{\medskip}

\newcommand{\bHuge}[1]{{\large {\bf #1}}}
\newcommand{\bLARGE}[1]{{ {#1}}}
\newcommand{\bLarge}[1]{{ {#1}}}
\newcommand{\blarge}[1]{{ {#1}}}

%\renewcommand{\newpage}{\bigskip}


\begin{document}

\bigskip

\begin{center}
\bHuge{Distribution laws 
for integrable eigenfunctions}


\bigskip\bigskip

\bLarge{Tatsuya Tate (Keio University)
\\ \bigskip
Joint work with\\ \bigskip
B.~Shiffman and S.~Zelditch
(Johns Hopkins University)}
\end{center}

\bigskip %%%%%%%%% PAGE 2  

\noindent \bHuge{$\bullet$ Plan of the Talk}

\bLARGE{

\begin{itemize}

\item[\S 1.] Background of problems on asymptotics of eigenfunctions

\begin{enumerate}
%\renewcommand{\labelenumi}{{\rm (\Alph{enumi})}}

\item[{\bf (A)}] What kind of problems?

\item[{\bf (B)}] For what kind of ``eigenfunctions''?

\end{enumerate}

\item[\S 2.] Toric manifolds and ``toric eigenfunctions''

(A short introduction to toric manifolds from a view point of integrable systems)

\item[\S 3.] Main results 

(``Logarithmic law'' for the limits of 
rescaled distribution functions for toric eigenfunctions)


\item[\S 4.] Explanation of the ``logarithmic law''. 
 
\end{itemize}

}


\bigskip %%%%%%%%%% PAGE 3


\begin{center}
\bHuge{\S 1. Background of problems on asymptotics of eigenfunctions}
\end{center}

%\medskip

\bLARGE{
We are going to study asymptotics of ``distribution function'' 
of ``eigenfunctions'' on a compact complex phase space. 

%\bigskip

\begin{enumerate}

\item What kind of problems on asymptotics of eigenfunctions?

%\bigskip

\item For what kind of eigenfunctions?

%\bigskip

\item Why such eigenfunctions?


\end{enumerate}

}


\bigskip


\noindent {\bf $\bullet$ Problems on Asymptotics of Eigenfunctions}

%\bigskip

\bLARGE{

\begin{itemize}

\item[(1)] Weak$^{*}$-limits of $|\varphi_{\lambda,\hbar}|^{2}d\vol\,$ or 
their microlocal lift as $\lambda \to \infty$ or $\hbar \to 0$ with $\lambda(\hbar) \in I$. 

%\bigskip

\item[(2)] Estimating the $L^{p}$-norms ($ 1 \leq p \leq \infty$) of eigenfunctions on configuration/phase space. 

%\bigskip

\item[(3)] Pointwise asymptotics. 

%\bigskip

\item[(4)] Asymptotics of distribution function 
\[
D_{\lambda}(t)=\vol(z\,;\,|\varphi (z)|^{2} >t). 
\]

\end{itemize}

}

%\bigskip

\bLARGE{

Our purpose is to give an answer to the problem (4) 
for a special kind of completely integrable systems on a phase space 
which is, in this talk, a compact K\"{a}hler toric manifold.

}

\bigskip 

\noindent {\bf $\bullet$ Pointwise asymptotics }

\bigskip

\bLARGE{Plancherel-Rotach formula for the Hermite functions: 
\[
\begin{gathered}
h_{n}(x)=c_{n}(-1)^{n}e^{x^{2}/2}
\left(
\frac{d}{dx}
\right)^{n}
(e^{-x^{2}}),\\
c_{n}=\frac{1}{\pi^{1/4}(2^{n}n!)^{1/2}},\qquad 
\|h_{n}\|_{L^{2}}=1. 
\end{gathered}
\]
}

\bigskip

\bLARGE{Set $\lambda_{n}:=2n+1$. Then $h_{n}$ is an eigenfunction of the 
harmonic oscillator $-\frac{d^{2}}{dx^{2}}+x^{2}$ with the 
eigenvalue $\lambda_{n}$.

}

\bigskip

\noindent {\bf Theorem: } \bLARGE{Let $0 < \varepsilon < \omega$, $0 < c$. 


\begin{enumerate}

\item[{\bf (a)}] For $x=\sqrt{\lambda_{n}}\cos \phi$, $\epsilon < \phi <\pi -\epsilon$, 
\[
\begin{split}
|h_{n}(x)|^{2}&=\frac{1}{\pi \sqrt{\lambda_{n}}\sin \phi} \times \\
&\left[
\sin^{2}\left( 
\frac{\lambda_{n}}{4}(\sin(2\phi)-2\phi)
+\frac{3\pi}{4}
\right)
+O(1/n),
\right]
\end{split}
\]



\item[{\bf (b)}] for $x=\sqrt{\lambda_{n}}\cosh \phi$, $\varepsilon \leq \phi \leq \omega$, 
\[
|h_{n}(x)|^{2}=
\frac{1}{2\pi \sqrt{\lambda_{n}}\sinh \phi}
e^{-\frac{1}{2}\lambda_{n}(\sinh \phi -2\phi)}
(1+O(1/n)),
\] 



\item[{\bf (c)}] for $x=\sqrt{\lambda_{n}}-2^{-1/2}3^{-1/6}t$, $|t| \leq c$, 
\[
|h_{n}(x)|^{2}=
\left(
\frac{6^{1/3}}{\pi}
\right)^{2}
\frac{1}{\lambda_{n}^{1/6}}
(A(t)^{2} +O(n^{-2/3})),
\]
where $A(t)$ is the Airy function. 

\end{enumerate}

}


\noindent{\bf $\bullet$ Real and Complex eigenfunctions}



\bLARGE{

\begin{itemize}

\item[(A)] Real eigenfunctions (Schr\"{o}dinger): 

 

Eigenfunctions of 
elliptic operators $\varphi_{\lambda}$ with eigenvalue $\lambda$ 
or $\varphi(\lambda(\hbar))$ with parameter $\hbar$ on a manifold $B$. 
Or their lift to the phase space $M=T^{*}B$: 

\begin{itemize}

\item[(1)] Their microlocal lift
\[
C^{\infty}(T^{*}M) \ni a \mapsto \ispa{\op_{\hbar}(a)\varphi_{\lambda},\varphi_{\lambda}} \in \mb{C}. 
\]

\item[(2)] On the phase space $\mb{R}^{2n}$, 
Wigner function 
\[
\begin{split}
(&W_{\hbar}\varphi_{\lambda,\hbar})(x,\xi) \\
&=(2 \pi \hbar)^{-n}\int_{\mb{R}^{n}}e^{-i\eta \xi/\hbar}\varphi_{\lambda}(x+\tf{1}{2} \eta)
\overline{\varphi_{\lambda,\hbar}(x-\tf{1}{2} \eta)}\,d\eta.
\end{split}
\]
or Husimi function
\[
(H_{\hbar}\varphi_{\lambda})(x,\xi)
=(W_{\hbar}\varphi * (\mbox{Gaussian}) )(x,\xi) \ (\geq 0).
\]

\end{itemize}



\item[(B)] Complex Eigenfunction (Bargmann): 
\[
\begin{array}{l}
\bullet\  \mbox{phase space}=\mb{C}^{m}, \\
\bullet\  \mbox{quantum state space}=H^{0}(\mb{C}^{m}) \\
=\{f: \mbox{holom. on }\mb{C}^{m}\,;\,
f \in L^{2}(\mb{C}^{m},e^{-|z|^{2}/2}dzd\bar{z})\}. 
\end{array}
\]



\noindent $\bullet$ Thus, the eigenfunctions on phase space should be in such a space 
as the Bargmann space. 

\medskip

\noindent $\bullet$ Our complex eigenfunctions $\varphi_{\lambda}$ 
are on a `compact' complex phase space $(M,\omega)$ 
with a holomorphic Hermitian line bundle 
\[
L \longrightarrow M \mbox{ with a Hermitian form }h 
\]
such that $c_{1}(L)=\omega$. 

\bigskip

\noindent $\bullet$ The eigenfunctions $\varphi_{\lambda} \in H^{0}(M,L^{\otimes N})$ 
($N=1/(\mbox{Planck const.} \in \mb{N}$) 
would be eigenfunctions of automorphism(s) on $H^{0}(M,L^{\otimes N})$. 

\medskip

\noindent $\bullet$ Analogue of the Husimi functions are of the form: 
$|\varphi_{N\lambda}(z)|^{2}_{h_{N}}$, $z \in M$
where $h_{N}$ is a Hermitian metric on $L^{\otimes N}$ 
induced by $h$. 

\medskip

\noindent $\bullet$ We are going to consider eigenfunctions for a completely integrable systems 
on a complex compact phase space. 

\end{itemize}

}



\noindent {\bf Why such eigenfunctions?}

\bigskip

\bLARGE{

Usually, problems on asymptotics of eigenfunctions are quite hard to resolve 
completely. 


\bigskip

$\Longrightarrow$ We need rather simple model for which one 
can completely determine the limit of distribution functions 
of eigenfunctions. 

\bigskip

$\Longrightarrow$ There are a lot of non-trivial examples 
of completely integrable systems on compact complex phase spaces, 
and they are rather easy to handle.

\bigskip

$\Longrightarrow$ Toric manifolds and the ``toric'' eigenfunctions (sections) on it. 


}

\newpage


\begin{center}
\bHuge{\S 2. Toric manifolds and ``toric eigenfunctions''}
\end{center}


\noindent \bLARGE{Let $(M,\omega)$ be a compact K\"{a}hler manifold 
(e.g., $\mb{C}P^{m}$). 
}

\bigskip

\noindent \bLARGE{$(M,\omega)$ is a toric manifold  

\smallskip

\noindent$\Longleftrightarrow$ $M$ is an algebraic variety with 
an action of the complex torus $(\mb{C}^{*})^{m}$ 
on $M$ which has a (Zariski) dense open orbit, and the $(\mb{C}^{*})^{m}$-action 
is an extension of the $(\mb{C}^{*})^{m}$ on $(\mb{C}^{*})^{m} \subset M$. 

\smallskip

\noindent $\sim$ A certain completely integrable system on $(M,\omega)$. 

}

\bigskip

\noindent \bLARGE{$\bullet$ Completely integrable system is: } 

\medskip

\bLARGE{
$f_{1},\ldots,f_{m} \in C^{\infty}(M)$ ($\dim_{\mb{R}} M=2m$) such that 

\begin{itemize}

\item $\{f_{i},f_{j}\}=0$, 

\item $df_{1},\ldots,df_{m}$ is linearly independent on (an open dense subset in) $M$. 

\end{itemize}

\noindent $\Longleftrightarrow$ We have a (Hamiltonian) $\mb{R}^{m}$-action on $M$: 
$$
\begin{gathered}
t \cdot z :=\varphi_{t_{1}}^{1} \comp \cdots \comp \varphi_{t_{m}}^{m}(z),
\qquad
t=(t_{1},\ldots,t_{m}) \in \mb{R}^{m},\qquad z \in M 
\end{gathered}
$$
with the ``moment map'' 
$$
\mu:M \to \mb{R}^{m},\qquad \mu(z)=(f_{1}(z),\ldots,f_{m}(z)). 
$$

\bigskip

\noindent $\bullet$ To obtain a toric manifold from the above 
completely integrable system on a compact K\"{a}hler manifold $(M^{2m},\omega)$, 
we need the following 

\bigskip

\noindent $\bullet$ Requirement: \hspace{3pt} The Hamiltonian $\mb{R}^{m}$-action 
reduces to a global (effective) $T^{m}$-action on $M^{2m}$

\bigskip

\noindent $\bullet$ Fact (Delzant, Guillemin). $(M,\omega)$ with the above requirement has 
a structure of the toric manifold with the ``moment polytope''
\[
P:=\mu(M) \subset \mb{R}^{m}
\]
which is a polytope (convex hull of a finite set). 

}

\bigskip

\noindent $\bullet$ Toric eigenfunctions are eigenfunctions of a completely integrable system 
expressed by a (Hamiltonian) real torus action ${\bf T}^{m}$ on a 
compact K\"{a}hler toric manifold 
\[
M_{P}=\mbox{phase space of real dimension } 2m
\]
with the moment polytope $P$. 

\bigskip

\noindent $\bullet$ These are nothing but monomials, 
and which are realized as sections 
of a line bundle over a compact toric manifold: 
$$
L_{P} \longrightarrow M_{P}. 
$$


\noindent {$\bullet$ Fact: }

\medskip

\bLARGE{

\noindent $\forall P \subset \mb{R}^{m}$ : integral Delzant polytope, \\
$\exists (M_{P},\omega_{P})$ : a compact K\"{a}hler toric manifold,  
$\dim_{\mb{R}}M_{P}=2m$, \\
$\exists L_{P} \to M_{P}$ : a hol. line bundle with $c_{1}(L_{P})=\omega_{P}$ s.t., \\
$\forall N (=1/\hbar) \in \mb{N}$, $H^{0}(M_{P},L_{P})$ 
is decomposed into the eigenspaces (weight spaces) of ${\bf T}^{m}$-action 
as follows: 
$$
H^{0}(M_{P},L_{P}^{\otimes N}) =
\bigoplus_{\gamma \in NP \cap \mb{Z}^{m}}
\mb{C}\cdot \chi_{\gamma}^{P},
$$
where the eigenfunctions $\chi_{\gamma}^{P}$ are essentially the monomials on the 
open dense set $(\mb{C}^{*})^{m} \subset M_{P}$: 
$$
\chi_{\gamma}^{P}(z)=z^{\gamma}=z_{1}^{\gamma_{1}} \cdots z_{m}^{\gamma_{m}},
\quad z=(z_{1},\ldots,z_{m}) \in (\mb{C}^{*})^{m}.
$$

\noindent $\bullet$ The monomials $\chi_{\gamma}^{P}$ ($\gamma \in NP \cap \mb{Z}^{m}$) are orthogonal 
w.r.t. 
$$
\ispa{s,t}_{N,P}=
\int_{M_{P}}
h_{N}(s(z),t(z))\,\omega_{P}^{m}/m!, 
$$
given by the Hermitian structure. 
But they are not orthonormal. 

\medskip

$\Longrightarrow$ The normalized eigenfunctions 
\[
\varphi_{\gamma}^{P}:=\chi_{\gamma}^{P}/\|\chi_{\gamma}^{P}\|_{N,P}
\]
form an orthonormal basis for $H^{0}(M_{P},L_{P}^{\otimes N})$. 

$\Longrightarrow$ call it ``toric'' eigenfunctions.
}

\newpage
\noindent {\bf $\bullet$ Basic Example}

\bigskip

\bLARGE{We shall show how to construct $M_{P}$ by an example. 
Let $m=1$, $P=[0,p]$ with $0 <p \in \mb{Z}$. 
Then, the monomial embedding
\[
\Phi_{P}:\mb{C}^{*} \ni z \mapsto 
[c_{0}:c_{1}z:\cdots:c_{p-1}z^{p-1}:c_{p}z^{p}] \in \mb{C}P^{p}
\]
is an embedding, where $c_{\alpha}={p \choose \alpha}^{1/2}$, 
the binomial coefficients. 
Then the Zariski closure 
\[
M_{P}=\overline{\Phi_{P}(\mb{C}^{*})}^{{\rm Zariski}}
\]
is the image of the Veronese map: 
$$
\Psi:\mb{C}P^{1}\ni [z_{0}:z_{1}] 
\mapsto [c_{0}z_{0}^{p}:c_{1}z_{0}^{p-1}z_{1}:\cdots : c_{p-1}z_{0}z_{1}^{p-1}:c_{p}z_{1}^{p}]
\in \mb{C}P^{p}.
$$
Thus, we have $M_{P}=\mb{C}P^{1}$. The torus ${\rm T}^{1}=S^{1}$-action is given by 
\[
e^{i\theta}[\zeta]=\Psi([z_{0}:e^{i\theta}z_{1}]),\quad [\zeta]=\Psi([z_{0}:z_{1}]) \in M_{P}. 
\]
It is not hard to see that the K\"{a}hler form $\omega_{P}=\Psi^{*}\omega_{{\rm FS}p}$, where 
$\omega_{{\rm FS}p}$ is the Fubini-Study K\"{a}hler form on $\mb{C}P^{p}$, is given by 
$\omega_{P}=p\omega_{{\rm FS}1}$, where $\omega_{{\rm FS}1}$ is the Fubini-Study K\"{a}hler form 
on $\mb{C}P^{1}$. Hence $L_{P}:=\Psi^{*}\mathcal{O}(1)=\mathcal{O}(p)$, the $p$-th tensor power 
of the hyperplane line bundle over $\mb{C}P^{1}$. 
The moment map $\mu_{P}:M_{P}=\mb{C}P^{1} \to \mb{R}$ is given by 
\[
\mu_{P}(\Psi[z_{0}:z_{1}])=
\frac{p|z_{1}|^{2}}{|z_{0}|^{2}+|z_{1}|^{2}}.
\]
The weight space decomposition for $H^{0}(M_{P},L_{P}^{\otimes N})=H^{0}(\mb{C}P^{1},\mathcal{O}(Np))$ 
is given by 
\[
\begin{gathered}
H^{0}(\mb{C}P^{1},\mathcal{O}(Np))=
\bigoplus_{\alpha=0}^{Np}\mb{C}\cdot \wh{\chi}_{\alpha},\\
\wh{\chi}_{\alpha}([z_{0}:z_{1}])=z_{0}^{Np-\alpha}z_{1}^{\alpha},\quad \alpha \in [0,Np] \cap \mb{Z}. 
\end{gathered}
\]
Thus, on the open orbit $\mb{C}^{*}=\{[1:z] \in \mb{C}P^{1}\,;\,z \in \mb{C}^{*}\}$, 
the weight vector $\wh{\chi}_{\alpha}$ is 
identical to $\chi_{\alpha}(z)=z^{\alpha}$, $z \in \mb{C}^{*}$. 





}

For $x \in P$ and $\alpha_{N} \in (NP) \cap \mb{Z}^{m}$, 
we call the sequence $\{\alpha_{N}\}_{N \geq 1}$ an {\it approximately 
multiple (AP)} of $x \in P$ if 
$$
\alpha_{N}=Nx +O(1). 
$$
(We regard the sequence $\{\alpha_{N}\}$ as a ``semiclassical ray'', 
which is, in a usual setting, $Nx$ when $x \in P \cap \mb{Z}^{m}$)


\bigskip

\noindent {\bf Our Problem is:}

\bigskip

\bLARGE{

To find the asymptotic formula for the distribution function 
\[
D_{\alpha_{N}}(t):=\vol_{M_{P}}(z \in M_{P}\,;\,|\varphi_{\alpha_{N}}^{P}(z)|_{N,P}^{2}>t),\quad t>0 
\]
for the toric eigenfunctions $\varphi_{\alpha_{N}}^{P}$. 

}

\bigskip

\noindent {\bf Related Works}

\bLARGE{

\begin{itemize}

\item Kurlberg-Rudnick (mathematics) 

They obtains the limit distribution of the so-called ``Hecke eigenfunctions'' 
of a quantum cat map on the $2$-dim torus. 
Their eigenfunctions are that on the phase space. 
Their result is that the value distribution tends to 
the Wigner's semi-circle measure. 

\item Mirlin et al (physics) 

They considers a disordered system with a localization. 
They considers, in particular, a ``tail'' of eigenfunctions on a configuration space. 
Their results are universal, and they obtains ``logarithmic law'', which is 
quite similar to one of our result below. 
However, we consider rather ``center'' of the localization of eigenfunctions. 

\end{itemize}

}
\bigskip

\begin{center}
\bHuge{\S 3. Main Theorems}
\end{center}


\bLARGE{

\noindent In what follows, we assume that $x \in P^{o}$.  

}

\bigskip



\noindent {$\bullet$ \bf Un-rescaled distribution}

\medskip

\bLARGE{

\begin{maintheo}
Let $\alpha_{N} \in NP \cap \mb{Z}^{m}$ be an asymptotic multiple of $x \in P^{o}$. 
Then, we have 
$$
D_{\alpha_{N}}(t) \sim 
\frac{(\pi m)^{m/2}}{c(P,x)\Gamma(m/2 +1)}
\left(
\frac{\log N}{N}
\right)^{m/2}, 
$$
which does not depend on $t$, where the constant $c(P,x)$ is 
given by 
the constant $c(P,x)$ is 
$$
c(P,x)=\lim_{N \to \infty}
\left( 
\tf{N}{2\pi}
\right)^{-m/2}
\|\varphi_{\alpha_{N}}^{P}\|_{\infty}^{2}.
$$
\end{maintheo}

}

\medskip

\bLARGE{

Un-rescaled distribution functions have the universal 
limit, i.e., geometric data appear only in 
the constant $c(P,x)$. 

}


\bigskip

\noindent {$\bullet$ \bf Rescaled distribution:
%\phantom{$\bullet$ }
``Logarithmic law''}

\bLARGE{

\begin{maintheo}
Define the rescaled distribution function
$$
F_{\alpha_{N}}(t)
=
\left (\tf{N}{2\pi}
\right)^{m/2}
\vol_{M_{P}}
\left(z \in M_{P}\,;\,|\varphi_{\alpha_{N}}^{P}(z)|^{2}_{N} > (\tf{N}{2\pi})^{m/2}t
\right)
$$
for $t > 0$. Then, we have 
$$
\lim_{N \to \infty}F_{\alpha_{N}}(t)
=\frac{1}{c(P,x)\Gamma(m/2 +1)}
\left(
\log \frac{c(P,x)}{t}
\right)_{+}^{m/2},
$$
where the constant $c(P,x)$ is 
$$
c(P,x)=\lim_{N \to \infty}
\left( 
\tf{N}{2\pi}
\right)^{-m/2}
\|\varphi_{\alpha_{N}}^{P}\|_{\infty}^{2}.
$$
and $(\log s)_{+}$ denotes $\log s$ for $s \geq 1$ and $0$ otherwise. 
\end{maintheo}

}



\bLARGE{
The rescaled distribution functions $F_{\alpha_{N}}(t)$ have also a universal limit, 
only what the constant $c(P,\alpha)$ and the power $m/2$ depend on 
the geometry of $M_{P}$. 

\bigskip




$\Longrightarrow$ Natural question: 

Do the limits of distribution functions have any geometrical information?

\bigskip

Answer is : YES. 

}

\bigskip

\noindent {\bf $\bullet$ Distribution functions with large length rescaling}

\bLARGE{

\begin{maintheo} We have: 
\[
\begin{split}
\lim_{N \to \infty}D_{\alpha_{N}}(e^{-Nt}) &=\vol_{M_{P}}
(z \in (\mb{C}^{*})^{m}\,;\,b_{x}^{P}(z) <\log t) \\
&= \int_{\{\rho \in \mb{R}^{m}\,;\, b_{x}^{P}(\rho) <t \}}
\det A(\rho)\,d\rho,
\end{split}
\]
where, $A(\rho)$ is a real positive $(m \times m)$-matrix valued smooth function such that 
\[
c(P,x)=\frac{1}{\sqrt{\det A(\rho_{x})}}
\]
with the vector $\rho_{x}=\bar{\mu}_{P}^{-1}(x) \in \mb{R}^{m}$.  
The function $b_{x}^{P}$ is a function on $(\mb{C}^{*})^{m}$  
invariant under ${\bf T}^{m}$-action satisfying 
\[
\partial_{\rho}b_{x}^{P}(\rho)=
\bar{\mu}_{P}(\rho)-x, \quad b_{x}^{P}(\rho_{x}) =0
\]
For the map $\bar{\mu}_{P}$, see below. 
\end{maintheo}

}


\bigskip


\noindent {\bf $\bullet$ Remarks}

\bLARGE{

\begin{itemize}

\item The function $b_{x}^{P}$ is defined on
$(\mb{C}^{*})^{m}/{\bf T}^{m} \cong \mb{R}^{m}$. 
The moment map induces the diffeomorphism: 
$$
\bar{\mu}_{P}:\mb{R}^{m} \longrightarrow P^{o}, 
$$
$b_{x}^{P}$ has a unique critical point $\rho_{x}=\bar{\mu}_{P}^{-1}(x)$.

\item The function $b_{x}^{P}$ can not be extended to whole of $M_{P}$. 

\item The matrix-valued function $A(\rho)$, $\rho \in \mb{R}^{m}$ is given by 
$$
A(\rho)=\partial_{\rho}\bar{\mu}_{P}(\rho),
$$
which can be shown to be positive definite for all $\rho \in \mb{R}^{m}$. 

\item The function $b_{x}^{P}$ looks like a quadratic form. In fact, we have
$$
b_{x}^{P}(\rho_{x}+u)
=\ispa{A(\rho_{x})u,u}/2 +R_{3}(u). 
$$

\end{itemize}

}


\begin{center}
\bHuge{\S 3. Explanation of the ``logarithmic~law''}
\end{center}


\noindent {\bf $\bullet$ Pointwise asymptotics}

\medskip

\noindent \bLARGE{Let $\alpha_{N} \in NP \cap \mb{Z}^{m}$ be an (AP) of $x \in P^{o}$. 
Then we have 
\[
|\varphi_{\alpha_{N}}^{P}(z)|_{N}^{2}=
c(P,x)\left(
\frac{N}{2\pi}
\right)^{m/2}
e^{-Nb_{x}^{P}(z)}(1+O(N^{-1}))
\]
locally uniformly on $(\mb{C}^{*})^{m} \subset M_{P}$, where
\[
c(P,x)=\frac{1}{\sqrt{\det A(\rho_{x})}}.
\]

\noindent $\bullet$ By the Taylor expansion for the function $b_{x}^{P}$, we have 
\[
\begin{gathered}
\left(
\frac{N}{2\pi}
\right)^{-m/2}
|\varphi_{\alpha_{N}}^{P}(z)|_{N}^{2}=
\frac{e^{-\ispa{A(\rho_{x})u,u}}/2}{\sqrt{\det A(\rho_{x})}}
(1+O(N^{-1/2})),\\[10pt]
z=e^{\rho_{x}+u/\sqrt{N}}
\end{gathered}
\]
on a ball of radius $O(N^{-1/2})$ uniformly in $|u| \leq c$.
}


\bigskip



\noindent {\bf $\bullet$ A Reason for the logarithmic law}

\medskip

\bLARGE{

Concerning the above theorem, we consider the following Gaussian function
\[
g(u):=\frac{e^{-\ispa{Au,u}/2}}{\sqrt{\det A}},\quad u \in \mb{R}^{m},
\]
where $A$ is a real positive $(m \times m)$-matrix. 
Also consider the measure $\nu_{A}$ on $\mb{R}^{m}$ defined by 
\[
\nu_{A}:=\frac{\det A}{(2\pi)^{m/2}}\,du
\]
The measure $\nu_{A}$ is not a finite measure, 
but we have 
\[
\|g\|_{L^{k}(\nu_{A})}
=\frac{1}
{k^{m/2}(\det A)^{(k-1)/2}},\quad k=1,2,\ldots.
\]
Then, a direct computation shows: 
\begin{gather*}
\nu_{A}(u \in \mb{R}^{m}\,;\,g(u)>t)=
\frac{1}{c\Gamma(m/2 +1)}
\left(
\log \frac{c}{t}
\right)^{m/2}\\[5pt]
c=\frac{1}{\sqrt{\det A}},\quad 0 \leq t \leq c, 
\end{gather*}
which is exactly the same form in one of our main theorems. 

\medskip


$\Longrightarrow$ Thus we can say that:

\begin{center}
Monomials look like Gaussian around the corresponding invariant torus, 
and which makes the limit distribution obey the ``logarithmic law.'' 
\end{center}

}





\end{document}





